\documentclass[12pt,a4paper]{article}

% Packages
\usepackage{amsmath, amsfonts, amssymb, amsthm, bm, bbm, comment} 
\usepackage{enumerate, graphicx, xcolor, tcolorbox, hyphenat, circledsteps}
\usepackage[shortlabels]{enumitem}
\usepackage{url, hyperref}
\hypersetup{
	colorlinks   = true,
	urlcolor     = cyan,
	linkcolor    = red,
	citecolor    = black
}

% Formatting
\usepackage[a4paper, noheadfoot, scale={0.8,0.82}]{geometry}
\usepackage[T1]{fontenc}
\newcommand{\thetitle}{Model Solutions}
\newcommand{\theauthor}{Luca Bollini}
\newcommand{\thedate}{November 2026}
\hyphenpenalty=10000
\parindent0pt

% Commands
\newcommand{\N}{\mathbb{N}}
\newcommand{\Z}{\mathbb{Z}}
\newcommand{\Q}{\mathbb{Q}}
\newcommand{\R}{\mathbb{R}}
\newcommand{\eps}{\varepsilon}
\newcommand{\ddt}[1]{\frac{d}{d#1}}
\newcommand{\ddttwo}[1]{\frac{d^2}{d#1^2}}
\newcommand{\ts}{\thinspace}

% Colours
\newtcolorbox{questionbox}{
	colframe=black, 
	colback=white, 
	boxrule=0.2mm,
	arc=0mm,
}


% Specific packages
\usepackage{stmaryrd}
\usepackage{tikz}

\newcommand{\contradiction}{%
\begin{tikzpicture}[rotate=45,x=0.5ex,y=0.5ex]
\draw[line width=.2ex] (0,2) -- (3,2) (0,1) -- (3,1) (1,3) -- (1,0) (2,3) -- (2,0);
\end{tikzpicture}
}

\begin{document}

\framebox[16.8cm]{{\sc University of Warwick}\hfill{\sc Mathematics Institute}}

\vspace{2mm}
{\sc  \theauthor}\hfill{\sc \thedate}

\vspace{3mm}

\begin{center}
{\large \textbf {\underline{\thetitle}}}
\end{center}

\vspace{3mm}

1.i) from Week 7 Problems 1.

\begin{questionbox}
    Let $(s_n)$ and $(t_n)$ be sequences of real numbers, and suppose that $(s_n)$ converges. Does $(t_n)$ converge if there exists $0<r<1$ such that for all $n$ sufficiently large,
    \[
    |t_{n+1}-t_{n}| \leq r^n|s_{n+1}-s_n|\text{?}
    \]
\end{questionbox}

First recall that a sequence of real numbers converges if and only if it is Cauchy. Using the Cauchy property of $(s_n)$, let $\eps>0$ and let $N\in\N$ be sufficiently large so that both
\[
|t_{n+1}-t_{n}| \leq r^n|s_{n+1}-s_n|
\]
and
\[
|s_{m}-s_n|<(1-r)\eps
\]
for all $m,n>N$. Now, by the triangle inequality, we have for any $m>n>N$ that
\begin{align*}
    |t_{m}-t_{n}| &\leq |t_{m}-t_{m-1}|+|t_{m-1}-t_{m-2}|+\cdots+|t_{n+1}-t_{n}|\\
    &\leq r^{m-1}|s_{m}-s_{m-1}|+r^{m-2}|s_{m-1}-s_{m-2}|+\cdots+r^n|s_{n+1}-s_{n}|\\
    &\leq (r^{m-1}+r^{m-2}+\cdots+r^n)(1-r)\eps.
\end{align*}
We now see that the right hand side is bounded by the tail of a geometric series. That is,
\begin{align*}
    (r^{m-1}+r^{m-2}+\cdots+r^n)\leq r^n\sum_{k=0}^\infty r^k = r^n\frac{1}{1-r}\leq \frac{1}{1-r}.
\end{align*}
Consequently,
\begin{align*}
    |t_{m}-t_{n}| &\leq (r^{m-1}+r^{m-2}+\cdots+r^n)(1-r)\eps\leq \frac{1}{1-r}(1-r)\eps = \eps
\end{align*}
for any $m>n>N$. Since the assumption $m>n$ can be made without any loss of generality, it follows that $(t_n)$ is Cauchy, and therefore converges.
\end{document}